%=======================================================================
%  APPENDIX A: Frame-coding codebook (final version)
%  Paste after \appendix in Final_Report.tex, replacing the placeholder
%  \section{Codebook and coding prompt}.
%  Requires: booktabs, enumitem (both already loaded in the preamble).
%=======================================================================
\section{Frame-coding codebook}
\label{app:codebook}

This is the instrument used for all coding. It was applied by three coders: by
me on the 300-item validation set, by \texttt{claude-opus-5} on all
\Ncoded{} items, and by GABRIEL on \texttt{gpt-4o} as a cross-check on the
coarse frame only. The text supplied to the model is this codebook compiled into
a system prompt, cached as a fixed prefix so that every item was coded against
identical instructions. The full codebook, the pilot draft that preceded it, the
worked-example coding guide and the exact system prompt are released with the
code.

\subsection{Unit of analysis}

\begin{table}[htbp]\centering
\small
\begin{tabular}{ll}
\toprule
Source & What is coded\\
\midrule
News (MediaCloud) & the headline only; body text is not consulted \cite{liu2019}\\
Reddit            & the submission title plus the first 100 words of self-text\\
Hacker News       & the story title plus the first 100 words of story text\\
\bottomrule
\end{tabular}
\end{table}

HTML entities and tags are removed before coding, so the human coder and the
model see identical characters.

\subsection{Step 0: scope}

The coder asks whether the text makes a claim, however indirect, about AI's
effect on human work. If it does not, the item is coded \textsc{tangential},
all four function codes are set to \textsc{none}, and coding stops. Items
excluded at this step include AI product announcements with no labour claim, job
advertisements, technical discussion of AI with no reference to work, labour
news with no AI component, and items where AI and work appear together only
incidentally.

\begin{quote}
\textbf{R0. Topic is not the same as frame.} An item being \emph{about} AI and
jobs does not mean it \emph{frames} AI's labour effect. Code the
characterisation the text builds, not the subject it mentions. If it mentions
both but characterises neither, it is \textsc{tangential}.
\end{quote}

R0 is the codebook's response to the finding that automated framing methods
often recover topic instead \cite{nicholls2021}. It is also why the tangential
rate is treated as a statement about retrieval precision rather than about coder
reliability: an item enters the corpus by containing an AI term and a labour
term, and some sources meet that test far more often without saying anything
about the relationship between the two. Observed rates were 45.8\% for Hacker
News, 27.7\% for Reddit and 25.7\% for MediaCloud.

\subsection{Step 1: the coarse frame}

Exactly one label is assigned.

\begin{table}[htbp]\centering
\small
\begin{tabular}{p{2.6cm}p{11.4cm}}
\toprule
Code & Definition\\
\midrule
\textsc{replacement}    & AI takes over work people used to do; people leave the task, job or occupation.\\
\textsc{augmentation}   & AI extends or assists human work; people stay involved.\\
\textsc{transformation} & AI restructures what work is or how it is organised, which is neither straightforward loss nor straightforward gain.\\
\textsc{mixed}          & Two or more of the above carry roughly equal weight in the same item.\\
\textsc{tangential}     & Did not pass Step 0.\\
\bottomrule
\end{tabular}
\end{table}

\begin{description}[leftmargin=2.6em,style=nextline,nosep]
  \item[R1] \textbf{Dominance.} Where one frame is clearly in front, through
        headline emphasis, the opening clause or evaluative wording, code that
        frame even if another is present. Use \textsc{mixed} only where the two
        are genuinely balanced. The pilot draft set 15\% \textsc{mixed} as a
        warning level; the observed rate was 2.7\%.
  \item[R2] \textbf{Questions carry their assumption.} A headline asked as a
        question, such as ``Will AI take your job?'', is coded to the frame the
        question assumes. Posing something as a question does not remove the
        frame.
  \item[R3] \textbf{Denial still invokes.} A text arguing \emph{against} a frame
        is coded to the frame it engages, with the moral evaluation field
        capturing the stance. RQ1 asks which frames structure the discussion,
        not which ones people endorse.
  \item[R4] \textbf{Sarcasm.} Code the frame the words invoke unless the sarcasm
        is unmistakable from the item alone. No thread context is available.
  \item[R5] \textbf{No outside knowledge.} Code only what the text says, using
        nothing known about the linked article, the author, or later events.
\end{description}

\subsection{Step 2: the four Entman functions}

One value per function. \textsc{none} is a legitimate and common answer, because
headlines often carry only one or two of the four functions.

\begin{table}[htbp]\centering
\small
\begin{tabular}{ll}
\toprule
Function & Values\\
\midrule
Problem definition       & \textsc{job\_loss}, \textsc{skill\_change}, \textsc{opportunity}, \textsc{inequality\_power}, \textsc{none}\\
Causal attribution       & \textsc{technology}, \textsc{firms}, \textsc{policy}, \textsc{workers}, \textsc{none}\\
Moral evaluation         & \textsc{negative}, \textsc{positive}, \textsc{ambivalent}, \textsc{neutral}, \textsc{none}\\
Treatment recommendation & \textsc{regulate}, \textsc{adapt}, \textsc{redistribute}, \textsc{accelerate}, \textsc{none}\\
\bottomrule
\end{tabular}
\end{table}

Two distinctions matter in application. Under moral evaluation,
\textsc{neutral} means the text evaluates and lands on neither good nor bad,
whereas \textsc{none} means there is no evaluative content at all. Under
treatment recommendation, \textsc{none} dominates the field: 81.6\% of
replacement-framed items and roughly four fifths of all items recommend nothing.
This is a substantive result, but it also drives the low reliability figure
discussed below, and a coder should not invent a recommendation in order to
avoid choosing \textsc{none}.

A confidence rating of 1 to 3 and a free-text note accompany each item. Human
and model confidence agree only weakly (Krippendorff's $\alpha = 0.324$,
ordinal), so confidence is reported descriptively and never used as a weight, a
filter or a reason to exclude an item.

\subsection{Reliability}

The threshold was fixed in the proposal before any coding took place: Cohen's
$\kappa \geq \kappaTarget{}$ on three or more of the four Entman functions,
with a fallback to a three-category typology at $\kappa \geq 0.55$ if that were
not met.

\begin{table}[htbp]\centering
\caption{Agreement between the human coder and the primary model on the
\Nval{}-item validation set. Three of four functions passed, so the full
four-function scheme was retained and no fallback was used.}
\label{tab:app-reliability}
\small
\begin{tabular}{lcccc}
\toprule
Variable & $\kappa$ (all 300) & $\kappa$ (excl.\ tangential) & Raw agreement & Difference\\
\midrule
Coarse frame             & 0.797 & 0.765 & 0.860 & 0.063\\
Problem definition       & 0.834 & 0.860 & 0.887 & 0.053\\
Causal attribution       & 0.792 & 0.749 & 0.857 & 0.065\\
Moral evaluation         & 0.803 & 0.794 & 0.853 & 0.050\\
Treatment recommendation & 0.591 & 0.641 & 0.877 & 0.286\\
\bottomrule
\end{tabular}
\end{table}

Treatment recommendation is the one value below the threshold, and the reason is
the distribution of answers rather than the care taken in coding. Its raw
agreement of 0.877 is the second highest of the five variables, but because
roughly four fifths of items recommend nothing, a large amount of agreement is
expected by chance and the correction is severe. The final column of
Table~\ref{tab:app-reliability} shows this directly: the difference between raw
agreement and $\kappa$ is 0.286 for treatment against 0.050 to 0.065 for every
other variable. Both figures are reported so that the coefficient is not read as
evidence of unreliable coding. Consistent with that explanation, treatment is
the only variable whose $\kappa$ rises when tangential items are removed.

Per-ecosystem $\kappa$ ranges from 0.745 to 0.876 for the coarse frame and the
first three functions, indicating no drift in coding between sources. Treatment
varies more, from 0.467 on Hacker News to 0.683 on MediaCloud, which follows the
same explanation, since the share of items recommending nothing itself differs
by source.

\subsection{Known limitations of the instrument}

\begin{enumerate}[nosep,leftmargin=1.6em]
  \item \textbf{One human coder.} Reliability here mixes the quality of the
        instrument with my own consistency. Re-reading my earlier codes surfaced
        two items I would have coded differently later. A second coder on 50 to
        100 items would separate the two effects.
  \item \textbf{Implicit framing is not captured.} Both the relevance rule and
        the coding operate on explicit wording, so a post that implies job loss
        without naming AI or work in recognisable terms never reaches the coder.
  \item \textbf{Text length differs by source.} News items have a median of 13
        words against 106 to 109 for the two forums, and the four functions are
        easier to detect when there is more text. The length-split check in
        Section~4.3 holds, but re-coding the forum items on their titles alone
        would be a cleaner test and was not run.
  \item \textbf{No thread context.} R4 requires sarcasm to be unmistakable from
        the item alone, which will miss irony that a reader of the full thread
        would catch.
\end{enumerate}
